\chapter{Conclusion}
\label{ch:conclusion}

\epigraph{The purpose of models is not to fit the data but to sharpen the questions.}{--- Samuel Karlin}

\lettrine[lines=2]{T}{his chapter} synthesizes the contributions of the RCC agent-based model with glucose sensing, discusses its limitations, and outlines directions for future work.

% ======================================================================
\section{Summary of Contributions}
\label{sec:contributions}

The central biological contribution of this work is the introduction of a continuous three-dimensional glucose concentration field overlaid on the discrete agent grid, where glucose is injected by blood vessels, diffuses according to Fick's law, and is consumed by all cell types at metabolically appropriate rates. Tumor cells consume glucose at approximately 200 times the normal rate consistent with the Warburg effect, creating depletion zones that suppress immune function and viability, while immune cells use glucose gradients for chemotactic navigation. This metabolic layer, largely absent from prior agent-based tumor microenvironment models, couples with 18 distinct agent types whose interactions are mediated by an effect system encoding local signaling, a 16-gene DNA system with codon-level translation for modeling clonal evolution and neoantigen generation, and sex-dependent treatment response modeled through estrogen, testosterone, and progesterone pathways across 17 effect channels including a biphasic estrogen response curve.

On the computational side, the migration from Mesa to Repast4Py restructures the model for MPI-based distributed execution while preserving all agent types and interaction logic. Parameter optimization is handled by an Optuna-based Bayesian pipeline searching the 87-parameter space, and a Streamlit user interface provides accessible simulation configuration and visualization. The integrated model was validated through a 160-simulation study comparing predictions against ARON-1 clinical trial data, with sex-stratified results demonstrating 65\% survival for simulated female patients under ICI+TKI combination therapy versus 55\% for males---consistent with clinical literature on sex-dependent immunotherapy response in RCC.

% ======================================================================
\section{Model Limitations}
\label{sec:limitations}

The current model operates on a single MPI rank, restricting tissue volume and agent count; distributing the glucose field across multiple ranks would require halo exchange protocols that, while supported in principle by Repast4Py's \texttt{SharedContext}, remain unimplemented for the coupled agent--field system. The metabolic subsystem models only glucose via linear diffusion, omitting the coupled dynamics of oxygen tension, lactate accumulation, and amino acid availability that govern immune polarization and tumor acidosis in biological tissue. Blood vessel agents inject glucose at fixed rates without responding to hemodynamic conditions or angiogenic signaling, representing a static snapshot of what is a dynamic, treatment-responsive vascular system---a limitation particularly relevant for TKI simulations.

Immune cells are initialized at random grid positions without modeling vascular entry or spatial distribution relative to the tumor margin, preventing the model from capturing immune-excluded or immune-desert phenotypes. T~cell receptor matching uses a discrete threshold rather than continuous binding affinities, which may underestimate immune response breadth against heterogeneous tumors. Finally, multi-rank configurations have not yet been validated against the single-rank baseline, and ensuring numerical equivalence across rank configurations must precede any scaling studies.

% ======================================================================
\section{Future Work}
\label{sec:future-work}

The most immediate priority is multi-rank MPI scaling through domain decomposition and halo exchange, enabling larger tissue volumes and cell counts. Extending the metabolic system to include oxygen, lactate, and amino acid fields with coupled reaction-diffusion terms would capture metabolic crosstalk governing immune polarization and tumor acidosis, while replacing the fixed vessel topology with dynamic VEGF-driven angiogenesis would enable simulation of vascular remodeling under anti-angiogenic therapy. Over the medium term, tracking clonal lineage relationships would illuminate how spatial competition and immune selection shape tumor heterogeneity, integrating patient-specific geometry from imaging data would move toward personalized simulation, and GPU porting of the diffusion solver via CuPy or JAX would make large-scale parameter sweeps practical. Longer-term directions include prospective integration with clinical trial workflows pending regulatory validation, adaptation to other solid tumors such as melanoma and NSCLC to test architectural generality, and systematic \emph{in silico} screening of combination therapy schedules at scales infeasible in clinical trials.

% ======================================================================
\section{Clinical Impact}
\label{sec:clinical-impact}

This work demonstrates that mechanistic agent-based models calibrated against trial data can generate clinically relevant predictions, moving beyond population-level averages toward stratified treatment modeling. The coupling of metabolic dynamics with immune system modeling provides a framework for investigating differential patient responses, while virtual patient cohorts generated by sampling the validated parameter space enable \emph{in silico} trial design---estimating statistical power, identifying responsive subgroups, and testing endpoint sensitivity before committing to costly clinical evaluation. Systematic parameter variation can also isolate individual biological variables' contributions to outcome, identifying candidate biomarkers that warrant prospective validation.

% ======================================================================
\section{Lessons Learned}
\label{sec:lessons-learned}

The most consequential insight concerns multi-seed validation. Early in development, individual runs were used to assess parameter changes, occasionally producing conclusions that did not hold under different seeds. The inherent stochasticity of agent initialization, mutation events, and probabilistic interactions means any single run is a sample from a surprisingly wide outcome distribution. The adoption of ensemble analysis---formalized in the 160-simulation study---was essential for distinguishing genuine parameter effects from noise and producing meaningful confidence intervals.

The decision to incorporate a continuous glucose field, rather than relying solely on agent-based rules, proved a significant source of predictive value. Purely agent-based approaches struggle to represent the spatially heterogeneous metabolic landscape that shapes immune viability. The glucose field introduces emergent spatial structure---depletion zones, nutrient gradients guiding migration---that arises naturally from coupling diffusion physics with agent behavior, and was instrumental in reproducing differential outcomes between well-perfused and poorly perfused tumor regions.

% ======================================================================
\section{Closing Remarks}
\label{sec:closing-remarks}

The four objectives stated in Chapter~\ref{ch:introduction}---porting to Repast4Py, adding glucose microenvironment sensing, validating against ARON-1 data, and optimizing parameters via Bayesian optimization---have each been addressed. The resulting model integrates metabolic, genetic, hormonal, and pharmacological subsystems into a single agent-based framework reproducing sex-stratified treatment outcomes consistent with clinical observation. While substantial limitations remain, the model establishes an extensible platform for the future directions outlined above and reflects the broader trajectory of computational oncology toward patient-level mechanistic prediction through coupled metabolic, genetic, and immunological dynamics.
